\documentclass{article}
\usepackage[utf8]{inputenc}
\usepackage{a4wide}

\begin{document}

\section*{maxSumCircular}

กำหนดอาร์เรย์วงกลมของเลขจำนวนเต็ม\texttt{nums}มาให้ชุดหนึ่ง

ให้หาอาร์เรย์ย่อย\texttt{(subarray)}ที่ผ\ul{\textbf{ลบวกรวมมีค่าสูงสุด}}อาร์เรย์วงกลมหมายถึงอาร์เรย์ที่ปลายท้ายสุดต่อกลับไปที่ส่วนต้นสุด

นั่นคือตัวก่อนหน้าของตัวแรกสุดคือตัวสุดท้าย\\



\hfill\break



\textbf{ตัวอย่าง}



\textbf{}



\begin{quote}

\textbf{input}



6



-3 1 2 -2 0 3



\textbf{output}



4



{subarrayที่ผลบวกรวมมีค่าสูงสุด =\texttt{{[}\ 1,\ 2,\ -2,\ 0,\ 3\ {]}}}

\end{quote}



\subsection*{Input}
บรรทัดที่ 1 ขนาดของอาร์เรย์\texttt{nums}



บรรทัดที่ 2 ตัวเลขจำนวนเต็ม เว้นด้วยช่องว่าง 1 ช่อง



\subsection*{Output}
{ให้หาอาร์เรย์ย่อย}\texttt{(subarray)}{ที่ผ}ลบวกรวมมีค่าสูงสุด\\



\end{document}
